\chapter{Konjunktiv}
\label{cha:Konjunktiv}

Der Konjunktiv ist ein Modus, der Irrealität, Wünsche, Höflichkeit oder indirekte Rede ausdrückt.

\section{Konjunktiv I – Indirekte Rede}
\label{sec:KonjunktivI}

Der Konjunktiv I wird verwendet, um fremde Aussagen neutral wiederzugeben.

\subsection{Bildung}
\begin{center}
\begin{tabular}{|l|c|c|c|}
\hline
\textbf{Person} & \textbf{sein} & \textbf{haben} & \textbf{werden} \\ \hline
ich & sei & habe & werde \\ \hline
du & seiest & habest & werdest \\ \hline
er/sie/es & sei & habe & werde \\ \hline
wir & seien & haben & werden \\ \hline
ihr & seiet & habet & werdet \\ \hline
sie/Sie & seien & haben & werden \\ \hline
\end{tabular}
\end{center}

\subsection{Verwendung}
\begin{itemize}
    \item \textbf{Indirekte Rede}: Er sagt, er \textbf{sei} krank. (Original: „Ich bin krank.")
    \item \textbf{Zeitungsberichte}: Der Minister \textbf{teilte} mit, die Verhandlungen \textbf{seien} erfolgreich.
    \item \textbf{Wissenschaftliche Berichte}: Die Studie \textbf{zeige}, dass...
\end{itemize}

\section{Konjunktiv II – Irreale Wünsche und Höflichkeit}
\label{sec:KonjunktivII}

Der Konjunktiv II drückt Irrealität, Wünsche oder Höflichkeit aus.

\subsection{Bildung}
\begin{itemize}
    \item \textbf{regelmäßige Verben}: Präteritum + -e
        \begin{itemize}
            \item kommen $\rightarrow$ käme
            \item gehen $\rightarrow$ ginge
            \item sehen $\rightarrow$ sähe
        \end{itemize}
    \item \textbf{unregelmäßige Verben}: Umlaut im Präteritum
        \begin{itemize}
            \item haben $\rightarrow$ hätte
            \item sein $\rightarrow$ wäre
            \item werden $\rightarrow$ würde
        \end{itemize}
    \item \textbf{Modalverben}: Siehe Tabelle
        \begin{itemize}
            \item können $\rightarrow$ könnte
            \item müssen $\rightarrow$ müsste
            \item dürfen $\rightarrow$ dürfte
        \end{itemize}
\end{itemize}

\subsection{Höflichkeitsform}
\begin{itemize}
    \item \textbf{Könnten Sie mir helfen?} (= Bitte helfen Sie mir.)
    \item \textbf{Würden Sie das bitte wiederholen?} (= Bitte wiederholen Sie das.)
    \item \textbf{Ich würde gern wissen, ob...} (= Ich möchte gern wissen, ob...)
\end{itemize}

\subsection{Ireale Wünsche}
\begin{itemize}
    \item \textbf{Wäre ich nur reich!} (= Ich bin nicht reich, aber ich wünschte, ich wäre es.)
    \item \textbf{Wenn ich nur Zeit hätte!} (= Ich habe keine Zeit.)
    \item \textbf{Es wäre schön, wenn...} (= Es ist nicht so.)
\end{itemize}

\subsection{Ireale Bedingungssätze}
\begin{itemize}
    \item \textbf{Wenn ich Zeit hätte, würde ich kommen.}
    \item \textbf{Wäre ich früher gekommen, hätte ich ihn getroffen.}
\end{itemize}

\section{Übungen zum Konjunktiv}
\label{sec:KonjunktivUebungen}

\subsection{Konjunktiv I}
\begin{enumerate}
    \item „Ich komme morgen." $\rightarrow$ indirekt \quad \textit{(Lösung: Er sagt, er komme morgen.)}
    \item „Wir haben Hunger." $\rightarrow$ indirekt \quad \textit{(Lösung: Sie sagen, sie hätten Hunger.)}
    \item „Ich kann schwimmen." $\rightarrow$ indirekt \quad \textit{(Lösung: Er sagt, er könne schwimmen.)}
    \item „Es ist kalt." $\rightarrow$ indirekt \quad \textit{(Lösung: Er sagt, es sei kalt.)}
    \item „Ich bin fertig." $\rightarrow$ indirekt \quad \textit{(Lösung: Sie sagt, sie sei fertig.)}
    \item „Wir werden warten." $\rightarrow$ indirekt \quad \textit{(Lösung: Sie sagen, sie würden warten.)}
    \item „Ich lese gern." $\rightarrow$ indirekt \quad \textit{(Lösung: Er sagt, er lese gern.)}
    \item „Die Aufgabe war leicht." $\rightarrow$ indirekt \quad \textit{(Lösung: Sie sagen, die Aufgabe sei leicht gewesen.)}
    \item „Ich habe gelernt." $\rightarrow$ indirekt \quad \textit{(Lösung: Er sagt, er habe gelernt.)}
    \item „Wir müssen arbeiten." $\rightarrow$ indirekt \quad \textit{(Lösung: Sie sagen, sie müssten arbeiten.)}
\end{enumerate}

\subsection{Konjunktiv II}
\begin{enumerate}
    \item Ich möchte wissen, ob du Zeit hast. $\rightarrow$ höflich \quad \textit{(Lösung: Ich würde gern wissen, ob du Zeit hättest.)}
    \item Ich bin nicht reich. $\rightarrow$ Wunsch \quad \textit{(Lösung: Ich wäre gern reich.)}
    \item Kannst du mir helfen? $\rightarrow$ höflicher \quad \textit{(Lösung: Könntest du mir helfen?)}
    \item Ich habe keine Zeit. $\rightarrow$ Wunsch \quad \textit{(Lösung: Ich hätte gern Zeit.)}
    \item Bin ich du? $\rightarrow$ Irreal \quad \textit{(Lösung: Wäre ich du,...)}
    \item Ich möchte Kaffee. $\rightarrow$ höflich \quad \textit{(Lösung: Hätten Sie gern Kaffee?)}
    \item Wenn ich Zeit \underline{\hspace{1cm}}, \underline{\hspace{1cm}} ich. (haben/reisen) \quad \textit{(Lösung: hätte - würde reisen)}
    \item Ich \underline{\hspace{1cm}} gern in Berlin. (sein) \quad \textit{(Lösung: wäre)}
    \item \underline{\hspace{1cm}} du mir sagen, wie spät es ist? (können) \quad \textit{(Lösung: Könntest)}
    \item Ich \underline{\hspace{1cm}} gern ein Auto. (haben) \quad \textit{(Lösung: hätte)}
\end{enumerate}

\chapterend